% 1405/5/4
% H. Behrouz 
% Email : Texlivegroup@gmail.com



% بعد از هر تغییر در این قسمت کلید Ctrl+S را بزنید و اجرا و مشاهده pdf فقط در فایل Slide.tex   می‌باشد.


% در این قسمت می‌توانید پکیج‌های مورد نیاز را اضافه کنید.
\usepackage{amsmath,amssymb,amsfonts}
\usepackage{mathtools}
\usepackage{xcolor} 
\usepackage{array}
\usepackage{multirow}
\usepackage{tikz} 
\usepackage{graphicx}
\usepackage{subfigure}
\usepackage{listings}
\usepackage{ragged2e}
\usepackage{mathtools}
\usepackage{float}
\usepackage{caption}
\usepackage{pifont}
\usepackage{tasks}
\usepackage{verbatim}





% انتخاب تم برای اسلاید 
\usetheme{Madrid}
%\usetheme{AnnArbor}
%\usetheme{CambridgeUS}
%\usetheme{Copenhagen}
%\usetheme{Ilmenau}
%\usetheme{Warsaw}


\usefonttheme[stillsansserifmath]{serif}


% به جای blue  در سطر ۳۱  می‌توانید هر یک از رنگ زیر را استفاده کنید
% black, blue, brown, cyan, darkgray, gray, green, lightgray, lime, magenta, olive, orange, pink, purple, red, teal, violet, white, yellow
\usecolortheme[named=blue]{structure}
\setbeamercovered{transparent}



% تعریف فونت 
\usepackage{xepersian}
\settextfont{Yas}
\setdigitfont{Yas}
\defpersianfont\nas[Scale=1.5]{IranNastaliq}


% دستور لوگو
\newcommand{\nologo}{\setbeamertemplate{logo}{logo}}


% سه خط دستور زیر در صورت فعال نبودن  شماره اسلاید فعال کنید
%\expandafter\def\expandafter\insertshorttitle\expandafter{%
%\insertshorttitle\hfill%
%\inserttotalframenumber\,/\,\insertframenumber}


% دستور ایجاد شماره برای  عکس و جدول
\setbeamertemplate{caption}[numbered]











